% =====================================================================
% APPENDIX A — GLOSSARY OF TERMS
% =====================================================================
\appendix
\chapter{Glossary of Terms}
\label{app:glossary}
\begingroup\small\setstretch{1.0}

This appendix explains every technical term used in the report, in the order a
newcomer is likely to meet it: first the task and the language phenomenon, then
the models and their training, then the evaluation metrics---with their
formulas---and finally the statistical machinery that turns scores into claims.
Each entry is written to be self-contained; no entry assumes another.

\vspace{4pt}
\noindent\textbf{Code and demo availability.} The complete pipeline, the executed
notebook and all artifacts (result tables, figures, run manifest and
requirements freeze) are maintained in the project's GitHub repository with
regular commits from both team members. The working demo---the fully executed
Kaggle notebook, including every output and figure in this report---is available
at \url{https://www.kaggle.com/code/anishdate/code-mixed-hate-speech-detection-pipeline-v2}.

\section{Task and Language Terms}
\label{app:task}

\textbf{Hate speech detection.} The classification task, given a post, of deciding
whether it expresses hate or attacks a person or group (label 1, \emph{hate}) or
does not (label 0, \emph{non-hate}). This study uses the binary form; datasets
that distinguish target groups or severity are outside its scope.

\textbf{Code-mixed language.} Text in which two or more languages alternate
within the same sentence or phrase, without typographic separation---from a single
borrowed word to a full clause switch; the ``mixing index'' entry below gives the
measure used here.

\textbf{Hinglish.} The Hindi--English code-mixed register dominant on Indian
social media: Hindi written in the Roman script (``romanised Hindi''), freely
interspersed with English words. Written Hindi normally uses the Devanagari
script; the Roman-script form has no standard spelling.

\textbf{Romanisation.} Writing a non-Latin script with Latin letters. Because it
is done by ear, one Hindi word acquires many spellings (\emph{hai}, \emph{hain},
\emph{he}, \emph{hn})---the main reason word-level vocabularies fragment and
character-level or subword models cope better.

\textbf{De-identification.} Removing personal information from text before
analysis: here, replacing URLs, email addresses, phone numbers and @handles with
placeholders, and collapsing letter elongation. Regex-based de-identification does
not catch plain-text names (Table~\ref{tab:limits}).

\textbf{PII.} Personally identifiable information: names, handles, phone numbers,
addresses and similar data that can identify a person.

\textbf{Data card.} A metadata file written next to a processed dataset,
recording its source, revision, licence, split sizes and preprocessing rules, so
that the data behind an experiment is auditable.

\textbf{Class imbalance.} A mismatch between the number of examples of each
class. PRISM is mildly imbalanced (53.5\% / 46.5\%), which is why macro-averaged
metrics are preferred over raw accuracy.

\section{Models and Training Terms}
\label{app:models}

\textbf{TF-IDF.} Term frequency--inverse document frequency. A way of turning text
into numbers: each feature counts how often a term occurs in a post (TF), down-weighted
by how common the term is across all posts (IDF), so that ubiquitous words carry
little signal.

\textbf{Character $n$-gram.} A substring of $n$ consecutive characters
(\emph{hain} contains the 3-grams \emph{hai} and \emph{ain}). Char $n$-grams
match misspelled and romanised variants through shared substrings and need no
dictionary---the reason the B0 baseline uses them with $n$ from 2 to 5.

\textbf{Linear SVM.} A linear support vector machine: it learns a weighted
combination of features and classifies by the sign of the score. It has no
probability output (hence no threshold to tune) and trains in seconds.

\textbf{Baseline.} A deliberately simple reference model that defines the score a
new method must beat. B0 (TF-IDF + SVM) plays this role: without it, a
transformer's three-point gain would have no costed reference point.

\textbf{Transformer.} The neural architecture underlying modern NLP: all tokens
are processed in parallel through stacked blocks of self-attention and
feed-forward layers, trained by gradient descent.

\textbf{Self-attention.} The mechanism by which every token computes a weighted
view of every other token, letting the model connect, say, a pronoun with the
noun it refers to several words away.

\textbf{Pretrained language model.} A transformer that has already been trained on
a large unlabelled corpus and therefore knows spelling, syntax and some world
knowledge before the task-specific training starts. HingRoBERTa, MuRIL and XLM-R
are three such models with different pretraining corpora.

\textbf{Fine-tuning.} Continuing to train a pretrained model on a small labelled
dataset (here: 17{,}704 posts, 3 epochs) so that its general language knowledge is
re-pointed at the task. All transformer results here are fine-tuned models.

\textbf{Classification head.} The small layer added on top of a pretrained
encoder that maps its output to the two class scores. It starts from random
weights and is trained together with (or on top of) the encoder.

\textbf{Epoch, batch, learning rate.} One \emph{epoch} is a pass over the
training data; training proceeds in \emph{batches} of 32 posts; the \emph{learning
rate} ($2\times10^{-5}$) is the optimiser's step size.

\textbf{Warmup and weight decay.} \emph{Warmup} (ratio 0.1) raises the learning
rate gradually over the first 10\% of steps, avoiding instability at the start;
\emph{weight decay} (0.01) penalises large weights, a mild regulariser.

\textbf{fp16 (mixed precision).} Storing and multiplying most network values in
16-bit instead of 32-bit floats, roughly halving GPU memory and speeding up
training. Small numerical differences remain, which is why seeds are described as
``reproducible in practice'' rather than bit-exact.

\textbf{Seed.} The number that starts all pseudo-random processes (weight
initialisation, shuffling, dropout). Re-running with seeds 42, 43, 44 measures
how much of the score is luck.

\textbf{Checkpoint.} A snapshot of the model saved during training. The pipeline
saves one per epoch and restores the one with the best \emph{validation} macro-F1,
not the final one.

\section{Evaluation Metrics}
\label{app:metrics}

All metrics are computed on a \textbf{confusion matrix}, the $2\times2$ table of
predicted vs.\ true classes. For the hate class: TP = hate posts predicted hate,
FP = non-hate posts predicted hate (false alarms), FN = hate posts predicted
non-hate (misses), TN = non-hate posts predicted non-hate.

\textbf{Precision} (per class) answers: ``of the posts the model called hate, how
many really were?''
\begin{equation}
\mathrm{Precision}_{c} = \frac{TP_{c}}{TP_{c} + FP_{c}}
\end{equation}

\textbf{Recall} (per class) answers: ``of the posts that really were hate, how
many did the model catch?''
\begin{equation}
\mathrm{Recall}_{c} = \frac{TP_{c}}{TP_{c} + FN_{c}}
\end{equation}

\textbf{F1} (per class) is the harmonic mean of the two---a class scores well on
F1 only if both precision and recall are high:
\begin{equation}
F_{1,c} = \frac{2\,\mathrm{Precision}_{c}\cdot\mathrm{Recall}_{c}}
               {\mathrm{Precision}_{c} + \mathrm{Recall}_{c}}
\end{equation}

\textbf{Macro-F1}---the primary metric of this report---is the unweighted average
of the per-class F1 over the $C$ classes ($C=2$ here):
\begin{equation}
\mFmac = \frac{1}{C}\sum_{c=1}^{C} F_{1,c}
\end{equation}
``Macro'' means each class counts equally. A model that is excellent on the large
class and useless on the small one gets a low macro-F1, which is exactly the
sensitivity wanted for hate speech: missing the hate class must cost as much as
failing on non-hate. This is why a 99\% accuracy on 99\% non-hate data would be
worthless, while macro-F1 would expose it as near zero.

\textbf{Macro precision / macro recall} apply the same unweighted averaging to
precision and recall. They are reported so that a macro-F1 change can be traced
to fewer false alarms (precision up) or fewer misses (recall up).

\textbf{MCC --- Matthews correlation coefficient.} A single-number summary of the
whole confusion matrix that is robust to imbalance, ranging from $-1$ (always
wrong) through 0 (random) to $+1$ (perfect):
\begin{equation}
\mathrm{MCC} =
\frac{TP\cdot TN - FP\cdot FN}
     {\sqrt{(TP+FP)(TP+FN)(TN+FP)(TN+FN)}}
\end{equation}

\textbf{Balanced accuracy} is the average of the two class recalls:
\begin{equation}
\mathrm{Balanced\ accuracy} = \tfrac{1}{2}\left(\frac{TP}{TP+FN} + \frac{TN}{TN+FP}\right)
\end{equation}
It ignores precision entirely, so comparing it with macro-F1 shows how much of a
model's quality comes from precision (fewer false alarms).

\textbf{Support.} The number of true examples of a class in the evaluated set.
Per-class scores are read against their support: the highly-mixed bucket of
Table~\ref{tab:buckets} carries only 72 posts, which is why its numbers move
widely.

\textbf{Decision threshold.} Transformers output a probability $p$ that a post is
hate; the \emph{threshold} $t$ turns it into a decision (hate if $p \geq t$).
The default $t = 0.5$ is a convention, not an optimum; this study tunes $t$ on
the validation set (\S\ref{sec:protocol}).

\clearpage
\section{Statistical and Protocol Terms}
\label{app:stats}

\textbf{Bootstrap confidence interval (CI).} A measure of how much a reported
number would wobble if the test set were redrawn. The test set is resampled with
replacement 1{,}000 times, the metric recomputed on each resample, and the 2.5th
and 97.5th percentiles form the 95\% CI. Within each resample the three seeds are
averaged, so the interval surrounds the plotted 3-seed mean.

\textbf{Seed variance / standard deviation.} The spread of a metric across seeds
42, 43, 44. It measures training luck, not test-set luck; both must be small
before two models can be ranked.

\textbf{McNemar's test.} A significance test for two models on the same items;
only the \emph{discordant} pairs matter---items where exactly one model is
correct. Under the null hypothesis ``both are equally good'', the split
$n_{10}$ vs.\ $n_{01}$ is a fair coin toss, and the exact binomial test on
$\min(n_{01}, n_{10})$ out of $n_{01}+n_{10}$ gives the $p$-value. It is the
correct paired test here because both models label the identical 8{,}852 posts.

\textbf{$p$-value.} The probability, assuming two models are in fact equally
good, of seeing a discrepancy at least as large as the observed one. A small
$p$-value is evidence \emph{against} equality, not a proof of importance.

\textbf{Holm correction.} When several hypotheses are tested at once (here: three
transformers each vs.\ the baseline), the chance of a false ``significant'' grows.
The Holm procedure orders the $p$-values and multiplies the $i$-th smallest by
$(m - i + 1)$, capping at 1, which keeps the family-wise error rate at the nominal
level. All ``significant'' claims in this report use Holm-adjusted $p$-values.

\textbf{Zero-shot transfer.} Applying a model to a domain or language it was never
trained on---here, the PRISM-trained model applied to Marathi--English posts with
no retraining and the same threshold. The 0.3770 result shows the term is not a
formality: zero-shot can fail completely.

\textbf{Data leakage.} Any route by which information from the test set influences
training or model selection, inflating reported scores. The v1$\rightarrow$v2 fix
(\S\ref{sec:protocol}) removed one concrete route: selecting the backbone on test
macro-F1.

\textbf{Validation set / test set.} The validation set (2{,}950 posts) is the
``exam for feedback'': all tuning and selection decisions read it. The test set
(8{,}852 posts) is the ``final exam'', touched once per seed for reporting.
Keeping the channels separate is what makes the reported numbers mean what they
claim.
\endgroup
